\begin{lemma}
\label{lemma-direct-image-coherator}
Let $f : X \to Y$ be a morphism of schemes.
Assume $X$ and $Y$ are quasi-compact and have affine diagonal.
Then, denoting
$$
\Phi : D(\QCoh(\mathcal{O}_X)) \to D(\QCoh(\mathcal{O}_Y))
$$
the right derived functor of
$f_* : \QCoh(\mathcal{O}_X) \to \QCoh(\mathcal{O}_Y)$
the diagram
$$
\xymatrix{
D(\QCoh(\mathcal{O}_X)) \ar[d]_\Phi \ar[r] &
D_\QCoh(\mathcal{O}_X) \ar[d]^{Rf_*} \\
D(\QCoh(\mathcal{O}_Y)) \ar[r] &
D_\QCoh(\mathcal{O}_Y)
}
$$
is commutative.
\end{lemma}

\begin{proof}
Observe that the horizontal arrows in the diagram are
equivalences of categories by
Proposition \ref{proposition-quasi-compact-affine-diagonal}.
Hence we can identify these categories (and similarly for
other quasi-compact schemes with affine diagonal).
The statement of the lemma is that the canonical map
$\Phi(K) \to Rf_*(K)$ is an isomorphism for all $K$ in
$D(\QCoh(\mathcal{O}_X))$. Note that if $K_1 \to K_2 \to K_3 \to K_1[1]$
is a distinguished triangle in $D(\QCoh(\mathcal{O}_X))$ and
the statement is true for two-out-of-three, then it is true
for the third.

\medskip\noindent
Let $U \subset X$ be an affine open. Since the diagonal of $X$ is affine,
the inclusion morphism $j : U \to X$
is affine (Morphisms, Lemma \ref{morphisms-lemma-affine-permanence}).
Similarly, the composition $g = f \circ j : U \to Y$ is affine.
Let $\mathcal{I}^\bullet$ be a K-injective complex in $\QCoh(\mathcal{O}_U)$.
Since $j_* : \QCoh(\mathcal{O}_U) \to \QCoh(\mathcal{O}_X)$
has an exact left adjoint
$j^* : \QCoh(\mathcal{O}_X) \to \QCoh(\mathcal{O}_U)$
we see that $j_*\mathcal{I}^\bullet$ is a K-injective complex
in $\QCoh(\mathcal{O}_X)$, see
Derived Categories, Lemma \ref{derived-lemma-adjoint-preserve-K-injectives}.
It follows that
$$
\Phi(j_*\mathcal{I}^\bullet) =
f_*j_*\mathcal{I}^\bullet =
g_*\mathcal{I}^\bullet
$$
By Lemma \ref{lemma-affine-pushforward} we see that
$j_*\mathcal{I}^\bullet$ represents $Rj_*\mathcal{I}^\bullet$ and
$g_*\mathcal{I}^\bullet$ represents $Rg_*\mathcal{I}^\bullet$.
On the other hand, we have $Rf_* \circ Rj_* = Rg_*$.
Hence $f_*j_*\mathcal{I}^\bullet$ represents $Rf_*(j_*\mathcal{I}^\bullet)$.
We conclude that the lemma is true for any complex
of the form $j_*\mathcal{G}^\bullet$ with $\mathcal{G}^\bullet$
a complex of quasi-coherent modules on $U$. (Note that if
$\mathcal{G}^\bullet \to \mathcal{I}^\bullet$ is a quasi-isomorphism,
then $j_*\mathcal{G}^\bullet \to j_*\mathcal{I}^\bullet$ is a
quasi-isomorphism as well since $j_*$ is an exact functor
on quasi-coherent modules.)

\medskip\noindent
Let $\mathcal{F}^\bullet$ be a complex of quasi-coherent
$\mathcal{O}_X$-modules. Let $T \subset X$ be a closed subset
such that the support of $\mathcal{F}^p$ is contained in $T$
for all $p$. We will use induction on the minimal number $n$
of affine opens $U_1, \ldots, U_n$ such that
$T \subset U_1 \cup \ldots \cup U_n$. The base case $n = 0$ is trivial.
If $n \geq 1$, then set $U = U_1$ and denote $j : U \to X$ the
open immersion as above. We consider the map of complexes
$c : \mathcal{F}^\bullet \to j_*j^*\mathcal{F}^\bullet$.
We obtain two short exact sequences of complexes:
$$
0 \to \Ker(c) \to \mathcal{F}^\bullet \to \Im(c) \to 0
$$
and
$$
0 \to \Im(c) \to j_*j^*\mathcal{F}^\bullet \to \Coker(c) \to 0
$$
The complexes $\Ker(c)$ and $\Coker(c)$ are supported
on $T \setminus U \subset U_2 \cup \ldots \cup U_n$ and the result
holds for them by induction. The result holds for
$j_*j^*\mathcal{F}^\bullet$ by the discussion in the preceding
paragraph. We conclude by looking at the distinguished triangles
associated to the short exact sequences and using the initial
remark of the proof.
\end{proof}